% =====================================================================
% PREÂMBULO: configuração técnica do documento.
%
% Normalmente você NÃO precisa editar este arquivo para escrever o TCC.
% Ele carrega os "pacotes", que são extensões que acrescentam recursos ao
% LaTeX (imagens, tabelas, cores, links...). A ordem importa: alguns
% pacotes precisam ser carregados antes de outros.
% =====================================================================

% Acentuação e hifenização corretas em português com o pdfLaTeX.
\usepackage[T1]{fontenc}
\usepackage[utf8]{inputenc}

% Fonte do texto (Times por padrão). Para trocar, edite o arquivo abaixo.
% =====================================================================
% FONTE TIPOGRÁFICA  (a letra do documento, não a "fonte" de uma figura)
%
% Mantenha somente uma das opções abaixo ativa: as três definem a família
% usada no texto inteiro. O padrão é Times, via pacote mathptmx.
% Confirme a fonte exigida pelo seu curso antes de alterar.
% =====================================================================
\usepackage{mathptmx}

% Alternativa sem serifa (família Helvetica, não a fonte Arial proprietária):
% \usepackage{helvet}
% \renewcommand{\familydefault}{\sfdefault}

% Alternativa Latin Modern (a fonte padrão moderna do LaTeX):
% \usepackage{lmodern}


\usepackage{courier}      % fonte de largura fixa, usada em \codigo{...}
\usepackage{indentfirst}  % recua também o primeiro parágrafo de cada seção
\usepackage{graphicx}     % \includegraphics
\usepackage{float}        % posicionamento de figuras e tabelas
\usepackage{microtype}    % melhora o espaçamento entre letras e palavras
\usepackage{booktabs}     % \toprule, \midrule, \bottomrule das tabelas
\usepackage{multirow,array}   % células que ocupam várias linhas
\usepackage{tabularx}     % tabelas com a coluna X, de largura automática
\usepackage[table,xcdraw]{xcolor}   % cores

% Liga cada referência às páginas em que ela foi citada, acrescentando
% "Citado na página 7." ao fim de cada entrada. É um apoio à revisão, e não
% um elemento previsto na norma (manual IFPI 2024, seção 11.3) -- por isso
% vem DESATIVADO. Se quiser usá-lo enquanto escreve, remova o % da linha e
% lembre de comentá-la de novo antes da entrega.
% Precisa vir ANTES de abntex2cite (carregado em configuracoes/citacoes).
% \usepackage[brazilian,hyperpageref]{backref}

\usepackage{simplecd}         % necessário para pos-textual/opcionais/capa-cd
\usepackage[final]{pdfpages}  % insere PDFs prontos (ficha catalográfica)

% Estilo institucional do IFPI (capa, folha de rosto, folha de aprovação).
\usepackage{abntex-ifpi/abntex-ifpi}

% Diagramas UML. Vem desativado porque é um arquivo grande e a maioria dos
% trabalhos não usa. Se for desenhar diagramas UML, remova o % da linha.
% \usepackage{abntex-ifpi/tikz-uml}

% Os dados e as cores precisam existir antes da configuração do PDF.
% =====================================================================
% COMECE AQUI.
%
% Regra de ouro: mude apenas o texto que está entre as chaves { }.
% Não apague o comando (a palavra que começa com \), nem as chaves.
%
%   \autor{Nome do aluno}
%   ^^^^^^ comando        ^ conteúdo: é só isto que você troca
%
% Estes dados são reaproveitados na capa, na folha de rosto, no resumo,
% nas notas de rodapé e nas propriedades do PDF. Preencher aqui basta;
% você não precisa repetir seu nome em nenhum outro arquivo.
% =====================================================================


% ---------------------------------------------------------------------
% 1. CURSO E CAMPUS  (obrigatório)
%
% Definidos uma única vez e reaproveitados em todos os lugares abaixo.
% Se você trocar o curso aqui, a capa, a folha de rosto e a nota de
% rodapé mudam juntas.
% ---------------------------------------------------------------------
\newcommand{\meucurso}{Análise e Desenvolvimento de Sistemas}
\newcommand{\meucampus}{Campus Pedro II}


% ---------------------------------------------------------------------
% 2. IDENTIFICAÇÃO DO TRABALHO  (obrigatório)
% ---------------------------------------------------------------------

% Título do artigo. Se houver subtítulo, separe-o por dois-pontos.
\titulo{Título do trabalho: subtítulo (se houver)}

% Seu nome completo, como deve aparecer na capa.
\autor{Nome do aluno}

% Cidade e estado que aparecem no rodapé da capa e da folha de rosto.
\local{Pedro II, Piauí}

% Ano de conclusão. Aparece logo abaixo da cidade.
\data{2026}

% Nome do curso para o abnTeX2 (usado na capa de CD, se você ativá-la).
\nomedocurso{\meucurso}


% ---------------------------------------------------------------------
% 3. INSTITUIÇÃO  (obrigatório)
%
% ATENÇÃO: aqui o \\ no fim da linha é obrigatório — ele força a quebra
% de linha no cabeçalho da capa. Mantenha um \\ ao fim de cada linha,
% menos na última. (No texto corrido dos capítulos, ao contrário, você
% NÃO deve usar \\ para quebrar linhas.)
%
% A terceira linha traz o nome do curso como ele aparece no diploma
% ("TECNOLOGIA EM ..."). Cursos de bacharelado ou licenciatura devem
% trocar esse prefixo, ou apagá-lo.
% ---------------------------------------------------------------------
\instituicao{%
  INSTITUTO FEDERAL DE EDUCAÇÃO, CIÊNCIA E TECNOLOGIA DO PIAUÍ\\
  \MakeUppercase{\meucampus}\\
  TECNOLOGIA EM \MakeUppercase{\meucurso}}


% ---------------------------------------------------------------------
% 4. TEXTO DA FOLHA DE ROSTO  (obrigatório)
%
% \tipotrabalho vai só nas propriedades do PDF (não aparece na página).
% \preambulo é o parágrafo que fica no lado direito da folha de rosto,
% explicando a que o trabalho se destina. Cuidado: este \preambulo não
% tem relação com o "preâmbulo" do LaTeX (o arquivo estrutura/preambulo
% .tex, que carrega os pacotes). São coisas diferentes com o mesmo nome.
% ---------------------------------------------------------------------
\tipotrabalho{Trabalho de Conclusão de Curso (Artigo Científico)}

\preambulo{Trabalho de Conclusão de Curso (artigo científico) apresentado
como exigência parcial para obtenção do diploma do Curso de \meucurso\ do
Instituto Federal de Educação, Ciência e Tecnologia do Piauí, \meucampus.}


% ---------------------------------------------------------------------
% 5. ORIENTAÇÃO  (orientador obrigatório; coorientador opcional)
%
% Escreva a titulação junto com o nome, por exemplo "Prof. Dr. Fulano".
% O rótulo "Orientador:" / "Coorientador:" é inserido automaticamente.
% ---------------------------------------------------------------------
\orientador{Prof. Dr. Nome do orientador}

% Sem coorientação? Deixe as chaves vazias: a linha some da folha de rosto.
\coorientador{}


% ---------------------------------------------------------------------
% 6. BANCA EXAMINADORA  (opcional)
%
% Estes nomes só aparecem se você ativar uma folha de aprovação em
% estrutura/pre-textual.tex. Enquanto ela estiver desativada, preencher
% aqui não muda nada no PDF.
%
%   \imprimirfolhadeaprovacao             -> orientador + membro um + dois
%   \imprimirfolhadeaprovacaoduascolunas  -> os quatro nomes, em 2 colunas
%
% Tem só dois avaliadores? Use a primeira versão e deixe \membrotres vazio.
% ---------------------------------------------------------------------
\membroum{Prof. Me. Nome do primeiro membro da banca}
\membrodois{Prof. Me. Nome do segundo membro da banca}
\membrotres{Prof. Dr. Nome do terceiro membro da banca}

% Instituição impressa sob cada assinatura da banca. Descomente e edite
% apenas se a banca for de outra instituição.
% \renewcommand{\instituicaodabanca}{Instituto Federal do Piauí (IFPI)}


% ---------------------------------------------------------------------
% 7. DATA DE APRESENTAÇÃO  (obrigatório)
%
% CUIDADO: esta data já aparece no PDF, no rodapé da página do abstract,
% mesmo com a folha de aprovação desativada. Se você não trocar, sai
% "Data de aprovação: Dia de mês de ano" no trabalho entregue.
% Exemplo de preenchimento: \dataapresentacao{20 de novembro de 2026}
% ---------------------------------------------------------------------
\dataapresentacao{Dia de mês de ano}


% ---------------------------------------------------------------------
% 8. NOTAS DE RODAPÉ DOS AUTORES  (obrigatório)
%
% Aparecem como nota de rodapé na primeira página, junto dos nomes.
% Repare que aqui há DOIS pares de chaves: o primeiro é o nome do
% comando, o segundo é o texto. Troque só o texto do segundo par.
%
%   \newcommand{\notadoautor}{este texto aqui}
%                ^nome^      ^--- o que você edita ---^
%
% Troque também os endereços: não deixe example.com no trabalho final.
% ---------------------------------------------------------------------
\newcommand{\notadoautor}{Graduando(a) em \meucurso\ pelo IFPI.
Contato: \url{aluno@example.com}.}

\newcommand{\notadoorientador}{Professor(a) do Instituto Federal do Piauí.
Contato: \url{orientador@example.com}.}


% ---------------------------------------------------------------------
% 9. PALAVRAS-CHAVE  (obrigatório)
%
% Texto simples, sem comandos de formatação, separado por ponto e vírgula
% e terminado em ponto final. \keywords é a versão em inglês, usada no
% abstract; traduza as mesmas palavras, não escolha outras.
% ---------------------------------------------------------------------
\newcommand{\palavraschave}{modelo; artigo; pesquisa.}
\newcommand{\keywords}{template; article; research.}

% =====================================================================
% CORES
%
% Os links do PDF (sumário, citações, referências cruzadas, endereços)
% são clicáveis mesmo quando impressos em preto. A cor abaixo vale para
% todos eles; quem define isso é configuracoes/pdf.tex.
% =====================================================================
\definecolor{cor-link}{RGB}{0,0,0}

% Para links azuis, comente a linha acima e ative esta:
% \definecolor{cor-link}{RGB}{8,40,75}

% Cor da borda dos quadros criados pelo ambiente "caixa" do estilo IFPI.
\definecolor{cor-quadro}{RGB}{5,28,63}

% PROPRIEDADES E LINKS DO PDF: dados definidos em metadados.tex.
\hypersetup{
  pdftitle={\imprimirtitulo},
  pdfauthor={\imprimirautor},
  pdfsubject={\imprimirtipotrabalho},
  pdfcreator={LaTeX com abnTeX2 e abntex-ifpi},
  pdfkeywords={\palavraschave},
  colorlinks=true,
  linkcolor=cor-link,
  citecolor=cor-link,
  filecolor=cor-link,
  urlcolor=cor-link,
  bookmarksdepth=4
}

% =====================================================================
% CITAÇÕES E REFERÊNCIAS
%
% No texto você usa \cite{chave} ou \citeonline{chave}; as obras ficam
% cadastradas em bibliografia.bib. Nunca escreva a lista de referências
% à mão: ela é montada a partir das citações.
%
% As opções abaixo são do pacote abntex2cite e definem a aparência:
%   alf ............ sistema autor-data, "(SILVA, 2020)". A alternativa
%                    seria o sistema numérico, "[1]".
%   versalete ...... imprime os sobrenomes em VERSALETE (letras maiúsculas
%                    de tamanho reduzido), como em Silva.
%   abnt-emphasize=bf .......... destaca o título em negrito
%   abnt-etal-list=3 ........... a partir de 3 autores, usa "et al."
%   abnt-etal-text=it .......... "et al." em itálico
%   abnt-and-type=& ............ liga os autores com "&"
%   abnt-last-names=abnt ....... formata os sobrenomes no padrão ABNT
%   abnt-repeated-author-omit=no  repete o nome do autor a cada entrada,
%                                 em vez de substituí-lo por um traço
%
% Confirme com seu curso qual sistema ele exige antes de mudar.
% =====================================================================
\usepackage[
  alf,
  versalete,
  abnt-emphasize=bf,
  abnt-etal-list=3,
  abnt-etal-text=it,
  abnt-and-type=&,
  abnt-last-names=abnt,
  abnt-repeated-author-omit=no
]{abntex2cite}

% Texto do link que leva da referência de volta à página em que foi citada.
\renewcommand{\backref}{}
\renewcommand*{\backrefalt}[4]{%
  \ifcase #1
    Nenhuma citação no texto.%
  \or
    Citado na página #2.%
  \else
    Citado nas páginas #2.%
  \fi}

% =====================================================================
% ESPAÇAMENTOS
% Manual de Normalização do IFPI (2024), seções 3.3, 3.4 e 7.4
%
% \parindent = recuo da primeira linha de cada parágrafo. A norma pede
%              1,25 cm a partir da margem esquerda (seção 3.4, alínea c).
% \parskip   = espaço extra entre parágrafos. A norma pede 0 pt: os
%              parágrafos são separados apenas pelo recuo.
%
% Não ajuste espaçamento parágrafo a parágrafo no texto: mude aqui e o
% documento inteiro acompanha.
% =====================================================================
\setlength{\parindent}{1.25cm}
\setlength{\parskip}{0pt}

% ---------------------------------------------------------------------
% ESPAÇAMENTO ENTRE LINHAS
%
% A regra geral (seção 3.3, alínea a) é 1,5 entre linhas. MAS o artigo
% científico é exceção: "fonte Arial ou Times New Roman, tamanho 12,
% espaçamento simples em todo o artigo" (seções 3.6.1 Nota 02 e 7.4).
%
% Este modelo é um ARTIGO, então o padrão aqui é o espaçamento simples.
% Se você adaptar o modelo para monografia, dissertação ou tese, troque
% \SingleSpacing por \OnehalfSpacing na linha abaixo.
%
% Em qualquer das duas modalidades continuam em espaço simples, por
% conta própria, as citações longas, as notas de rodapé, as referências,
% as legendas e a natureza do trabalho (seção 3.3, alínea b) -- isso o
% abnTeX2 já faz sozinho.
% ---------------------------------------------------------------------
\SingleSpacing


% Índice remissivo: ative esta linha somente se também ativar
% pos-textual/opcionais/indices.tex. Veja docs/guia-iniciante.md.
% \makeindex
